\section{Deciding which readings can be believed}

Between loading the CSVs and working out any chemistry sits one class, \texttt{Preprocessor}.
It takes the list of \texttt{Run} objects the loader built and returns a shorter list, with
some values inside the survivors set to \texttt{None}. It saves nothing and prints nothing.
Everything it removes is written to the \texttt{Journal} it was handed, so the report can say
later exactly what the modelling never saw.

Two things it deliberately does not do. It does not repair a reading: there is no filling in,
no smoothing, no interpolation. A value it cannot believe is removed, never replaced. And it
does not check that required fields are present. That is the loader's job and the loader is
all-or-nothing: each CSV is validated against its pandera schema as it is handed out, and a
table with a required field missing raises rather than loading. So by the time the
preprocessing sees a run, the run is well formed and only its values are in question. There is
a comment in \texttt{Preprocessor.steps} saying exactly this, marking where a fourth step used
to be needed and no longer is.

\subsection{The three steps, in order}

\texttt{Preprocessor.steps} is a tuple of three \texttt{Step} objects. Each one is given the
runs the step before it returned, so the order is the order of the pipeline.

\begin{enumerate}[nosep]
  \item \textbf{\texttt{DropExcludedRuns}} --- keeps only runs with \texttt{include} true.
        The mark is a decision taken at the bench, and this step only obeys it. If it would
        return nothing while it was given something, it raises \texttt{ValueError} instead:
        an empty result means the wrong file was loaded, not that every run was bad. Its
        journal label is \emph{not marked for use}.
  \item \textbf{\texttt{SkipSuspectReadings}} --- drops no run. It clears single values it
        cannot believe, and writes down why. Its label is \emph{suspect reading}.
  \item \textbf{\texttt{DropRunsMissingAnEndSample}} --- keeps only runs the lab has answered
        for: both lithium and magnesium, in both the concentrate and the dilute, at the
        first sampled minute and the last. Its label is \emph{no lab results at both ends}.
\end{enumerate}

Why step 3 comes last rather than first: it is cheaper to drop a run before inspecting its
readings, but then a run on its way out would never be looked over. Running the inspection
first means a doomed run still gets its bad values journalled, which is how the bench learns
that a particular sheet has a weight written down wrong even though that run never reached the
model. The real output shows this working --- run \texttt{AU} is dropped, and its rinse weight
at minute 240 is journalled anyway.

Step 3 asks only for the two ends because that is all the chemistry needs: the amount that
crossed the membrane is a difference between a start and a finish. The samples in between are
not asked for. The rinse is not asked for at all, because it is not part of the chemistry the
modelling uses. Both metals are required rather than lithium alone, because
permselectivity is a ratio of ratios and one metal cannot give it.

\subsection{Steps hold inspectors, inspectors hold checks}

Step 2 is split into three levels, and each level knows one thing.

\begin{description}[nosep,style=nextline]
  \item[A \texttt{Check} is a rule, and nothing else.] It is handed a list of readings and
        answers with the one reading at fault, which field to clear, and a sentence saying
        why --- or \texttt{None}. It never writes to the journal and never changes a reading.
        It is told at construction which field it judges and what to call it, so the same rule
        serves pH, temperature, weight or current without being rewritten.
  \item[An \texttt{Inspector} owns one part of a run.] It picks that part's readings out of the
        run, sorts them oldest first, builds the checks that part is held to, and does the
        clearing and the journalling for all of them.
  \item[The \texttt{Step} owns the loop over runs.] \texttt{SkipSuspectReadings} does nothing
        but put every run past every inspector.
\end{description}

The split is there so that a rule is written once and judged in one place. ``This value is
nothing like the rest of its own series'' is the same arithmetic for a tank weight and for the
supply current, and \texttt{ValueFarFromTheMedian} is that arithmetic, used twice with
different fields and different tolerances. It also keeps the thresholds out of the rules: no
check holds a number of its own, so a test can hand a check different numbers and watch it
behave differently without touching it.

There are four inspectors:

\begin{itemize}[nosep]
  \item \texttt{ConcTankInspector}
  \item \texttt{DiluteTankInspector}
  \item \texttt{RinseTankInspector}
  \item \texttt{PowerLogInspector}
\end{itemize}

\noindent The three tank ones share a parent,
\texttt{TankInspector}, which holds the construction of every check a tank may be given, so a
tank's own class is nothing but the list of what that tank is held to. As the code stands all
three tank lists are identical --- the same six checks in the same order. They are three
classes rather than one because the rules for one tank may well diverge from another's: the
dilute is topped up during a run and the concentrate is not, and the day a weight rule has to
differ between them, the change is a line in one class instead of an \texttt{if} inside a
shared one.

Each tank is also judged only against its own readings. A conductivity that is ordinary in the
dilute is impossible in the rinse, so there is no cross-tank comparison anywhere.

\subsection{What each check looks for}

Six check classes do all the judging. Three of them compare a reading with its own series,
two compare two instruments reading the same quantity, and one needs no comparison at all.

\subsubsection*{\texttt{ValueOutsideTheRange}}

Built for a value that is wrong on its own, whatever its neighbours say. It takes a low and a
high bound and rejects the first reading outside them:
\begin{equation}
  \label{eq:outside-range}
  \text{reject } x_i \quad\text{when}\quad \neg\,(\,\ell \le x_i \le h\,).
\end{equation}
It needs no minimum number of readings, because the judgement does not depend on any other
reading. It returns the \emph{earliest} offender rather than the biggest; order does not
matter here, because the inspector asks again after each one and works through them all. It
guards pH and temperature in every tank.

\subsubsection*{\texttt{ValueFarFromTheMedian}}

Built for a value that holds one level all run. With $m$ the median of the series and $s$ the
share allowed, it returns the reading furthest from $m$, if that distance is too far:
\begin{equation}
  \label{eq:far-from-median}
  i^\star = \arg\max_i\,\lvert x_i - m\rvert,
  \qquad
  \text{reject } x_{i^\star} \text{ when } \lvert x_{i^\star} - m\rvert > s\,\lvert m\rvert.
\end{equation}
The median is used rather than the mean because one bad value cannot drag it, and because it
does not care whether the bad value sits at the start of the run or the end. The tolerance is a
share of the median, not an absolute amount, so one number works for a rinse tank holding some
$4\unit{kg}$ and a concentrate tank holding some $2.1\unit{kg}$. It guards the tank weights and
the supply current, and it is the only check given the ``fewest readings'' threshold: with
fewer than three readings it says nothing, because with two each looks as odd as the other.

\subsubsection*{\texttt{ValueMovingTooFast}}

Built for a value that drifts but cannot lurch. Readings in this data are 5 to 240 minutes
apart, so the allowance is a rate and not a step --- strict on readings taken close together,
forgiving of slow drift across a long gap. For every consecutive pair it compares the speed of
the change with the allowance $a$ over $T$ minutes:
\begin{equation}
  \label{eq:moving-too-fast}
  \frac{\lvert x_{i+1} - x_i\rvert}{t_{i+1} - t_i} > \frac{a}{T}.
\end{equation}
Pairs no time apart are skipped, so the division is always safe. The fastest offending pair
wins. A pair says that two readings cannot both be right but not which one is wrong, so the
blame goes to whichever of the two sits further from the median of the series. It guards pH and
temperature in every tank, and needs at least two readings.

\subsubsection*{\texttt{JumpThatComesBack}}

Built for conductivity, which climbs or falls all run. There is no one level to compare
against, so the median is useless as a reference --- but a digit written down wrong still has a
signature: the series leaves its track and the next reading puts it straight back. A real
change of level stays changed. Over every window of three readings, with
$J = x_{i+1} - x_i$ and $B = x_{i+2} - x_{i+1}$,
\begin{equation}
  \label{eq:jump-and-back}
  \lvert J \rvert > g\,\lvert m \rvert
  \qquad\text{and}\qquad
  \frac{-B}{J} \ge c,
\end{equation}
and the middle reading $x_{i+1}$ is the one cleared. The first condition keeps an ordinary
wobble on a smooth series from counting as a jump; the median $m$ is used only to set that
scale. The second condition is the share of the jump that came back, and it is reported as a
percentage in the journal line. It can exceed 100\%: run \texttt{AR}'s concentrate at minute 60
put 132\% back, meaning the series overshot its way home. Fewer than three readings and the
check says nothing, since the window needs three. It is the one check with its own hard-coded
minimum rather than the shared threshold.

\subsubsection*{\texttt{PairedValuesTooFarApart}}

The stack voltage falls all run, so again there is no level to compare against --- but it was
written down twice, once from the supply and once from the multimeter. Those two agree to
hundredths of a volt, so a gap of a whole volt means one of them has a digit too many or too
few. The check takes the widest gap and rejects it if it is over the limit:
\begin{equation}
  \label{eq:voltage-gap}
  \text{reject the widest } \lvert v^{\text{supply}}_i - v^{\text{meter}}_i\rvert
  \quad\text{when it exceeds } \Delta_{\max}.
\end{equation}

\subsubsection*{\texttt{PairedValuesOutOfProportion}}

The same idea for the currents, which cannot be subtracted: the supply reads about a seventh
higher than the multimeter throughout, and in different units. So the two are held against
each other as a ratio, after the supply's amps are scaled by 1000 to meet the multimeter's
milliamps:
\begin{equation}
  \label{eq:current-ratio}
  r_i = \frac{1000\,I^{\text{supply}}_i}{I^{\text{meter}}_i}.
\end{equation}
The reading rejected is the one whose $r_i$ falls furthest from 1 among those outside
$[\,\ell,\,h\,]$. Readings where the multimeter wrote a zero are left out, so the ratio is
never divided by zero.

\subsubsection*{Which of two instruments gets the blame}

Both paired checks share \texttt{PairedValueCheck}, and both face the same problem: a
disagreement does not say which instrument is wrong. What does say is each column's own
history. The honest instrument stays near what it has been reading all run, so the blame goes
to whichever of the two sits further from the median of its \emph{own} column, measured as a
share of that median so that amps and milliamps compare fairly:
\begin{equation}
  \label{eq:odd-one}
  \text{blame} = \arg\max_{k \in \{\text{supply},\,\text{meter}\}}
  \frac{\lvert v^k_i - \operatorname{median}(v^k) \rvert}{\lvert \operatorname{median}(v^k)\rvert}.
\end{equation}
This is why the paired checks are the only ones that can clear a field other than the one they
are named after, and why their journal lines name the instrument. Run \texttt{AG} at minute 225
read $100.06\unit{V}$ on the supply against $10.00\unit{V}$ on the meter: a gap of
$90.06\unit{V}$, and the supply was the one nowhere near its own usual $11\unit{V}$, so the
supply's value was cleared. At minute 95 of the same run the supply read $11.20\unit{V}$ and the
meter $1.16\unit{V}$; this time the meter was the stranger in its own column, so the meter's
value went instead. Same check, same run, opposite verdicts, decided entirely by which column
the reading was out of step with.

For both paired checks a reading counts as judgeable only when \emph{both} instruments were
written down. One instrument alone is a reading with nothing to check it against, so it is
left alone.

\subsection{The numbers the checks judge by}

Every number a check compares against lives in one frozen dataclass called
\texttt{Thresholds}, with the reason for each one in the comment beside it. Change a number
there and every check that uses it changes with it. Table~\ref{tab:thresholds} is that file.

\begin{table}[htbp]
  \centering
  \caption{Every value in \texttt{thresholds.py}, and why it is what it is.}
  \label{tab:thresholds}
  \footnotesize
  \begin{tabular}{@{}llp{0.48\textwidth}@{}}
    \toprule
    Name & Value & Why \\
    \midrule
    \texttt{fewest\_readings\_to\_judge} & 3
      & Below this a check says nothing. With two readings each looks as odd as the other. \\
    \texttt{most\_of\_a\_series\_removable} & 0.34
      & The most of one series a check may clear. Needing more than a third means the file
        is wrong, not the readings. Past half, the median starts following the bad values. \\
    \texttt{impossible\_weight\_share} & 0.10
      & How far a tank weight may sit from its tank's usual weight. The three known bad
        weights are 90\% or more away; every good one is within 5\%. \\
    \texttt{most\_ph\_change} & 1.5
      & How far pH may move \ldots \\
    \texttt{most\_ph\_change\_in\_minutes} & 15.0
      & \ldots\ and over how long. A rate, because readings are 5 to 240 minutes apart. \\
    \texttt{plausible\_ph} & $(2.0,\,12.0)$
      & The range a pH reading must sit in whatever its neighbours say. These solutions cannot
        be at pH 0.07. \\
    \texttt{most\_temp\_change} & 2.0
      & How far temperature may move \ldots \\
    \texttt{most\_temp\_change\_in\_minutes} & 15.0
      & \ldots\ and over how long. The fastest real change in the data is $0.9\unit{^\circ C}$
        per quarter hour, so $2\unit{^\circ C}$ is a wide allowance. \\
    \texttt{plausible\_temp\_c} & $(5.0,\,40.0)$
      & The water sits at room temperature: every reading in the data is between 19.9 and
        $23.0\unit{^\circ C}$. \\
    \texttt{smallest\_jump\_share} & 0.10
      & How big a conductivity jump must be to count, as a share of what the series usually
        reads. Below a tenth, an ordinary wobble looks like a jump. \\
    \texttt{comes\_back\_share} & 0.7
      & How much of the jump the next reading must put back. At 0.7 this finds seven readings,
        every one a digit too many or too few. \\
    \texttt{impossible\_current\_share} & 0.10
      & How far the supply current may sit from its run's usual reading. It is held to
        within 2.2\% in nineteen of twenty-one runs. \\
    \texttt{largest\_voltage\_gap\_v} & 1.0
      & How far the two voltages may be apart. They agree to about $0.03\unit{V}$, so a gap
        this wide means a digit was written down wrong. \\
    \texttt{plausible\_current\_ratio} & $(0.62,\,1.6)$
      & The band the supply-to-meter current ratio must sit in. The supply reads about a
        seventh higher throughout, so a ratio and not a subtraction. \\
    \bottomrule
  \end{tabular}
\end{table}

Two of these are not any one check's. The fewest-readings number is handed only to the two
\texttt{ValueFarFromTheMedian} checks. The removable share belongs to the inspector, and no
check ever sees it.

\subsection{The one line where a value is cleared}

In the whole of the preprocessing there is exactly one place a reading's value is changed, in
\texttt{Inspector.\_clear\_what\_it\_rejects}:

\begin{quote}
\begin{verbatim}
reading, field_name, reason = worst
journal.skip(reading, field_name, reason)
setattr(reading, field_name, None)
\end{verbatim}
\end{quote}

That is the point of the design. The journal write is the line before the clearing, in the same
three lines, so there is no path through the code that clears a value without recording it.
Checks cannot clear anything --- they only answer with what they found --- so there is nowhere
else a value could quietly go missing. The journal entry takes its id and its record type from
the record itself rather than from the call site, so the id in the report and the reading that
was changed can never disagree.

\subsection{Asking a check again after every value it rejects}

A check is not run once over a series. It is asked for the single worst offender it can see;
if it finds one, that value is journalled and cleared, and the check is asked again from the
top. Two things make that the right loop.

The first is simply that a check answers with one reading, not a list, so a series holding two
bad values needs two passes. Run \texttt{AS}'s supply current held two: $0.46\unit{A}$ at
minute 210 and $0.44\unit{A}$ at minute 480, against a level of $0.045\unit{A}$. One pass would
have caught one of them.

The second is that the level a check compares against is recomputed each time, so it is a
better level once a bad value is gone. Every median, rate and ratio is worked out from the
values still present, which is what makes it safe to judge the second-worst reading against a
series the worst one has already left. How much it matters depends on the series. On
\texttt{AS}'s current the median read $0.045\unit{A}$ on both passes, because 57 readings do
not care about two of them. On run \texttt{AG}, where three voltage values went, the supply's
median moved only from $11.215$ to $11.19\unit{V}$ --- but the blame did move: the first pass
cleared the supply's $100.06\unit{V}$, and the two passes after it cleared the multimeter
instead.

The loop needs a stop. \texttt{Inspector.\_find\_clearing\_limit} sets how many times a single
check may fire on a single series:
\begin{equation}
  \label{eq:clearing-limit}
  \text{limit} = \max\!\left(1,\ \left\lfloor
    n \times \texttt{most\_of\_a\_series\_removable} \right\rfloor \right),
\end{equation}
where $n$ is the number of judgeable readings in the series \emph{before} any clearing. With
0.34, a pH series of three readings allows one removal, a typical tank series of seventeen
allows five, and a typical power log of forty-one rows allows thirteen. The $\max(1, \cdot)$
means a check can always fire at least once, even on the shortest series it will judge.

Two honest consequences. The first is the reason for the number: beyond about half a series the
median stops describing the good readings and starts describing the bad ones, at which point
the check turns around and clears the readings that were right. Stopping at a third keeps well
clear of that. If a third is not enough, the right conclusion is that the file is wrong, not
that the readings are.

The second is that \textbf{hitting the cap is silent}. The loop is a plain
\texttt{for \_ in range(limit)}; when it runs out of iterations it simply ends. Nothing is
journalled, nothing is printed, and the report will not tell you that a check still had
complaints when it was cut off. On the data as it stands this never happens: re-running the
step with the share raised to 0.5 and to 0.9 clears the same 27 values, so no check is being
held back. But that is a property of today's data, not a guarantee, and it is the one place in
this stage where something is left out without a line in the journal.

\subsection{Dropping a run against skipping a value}

The journal has two verbs, and they mean different things.

A \textbf{drop} is about a whole run. The run leaves the list and nothing downstream sees it.
A drop is written by \texttt{Preprocessor} itself, not by the step: after each step it compares
the run codes it handed in with the run codes it got back, and journals every run that did not
come back, with the step's own name as the reason. That is why a drop line reads
\emph{no lab results at both ends} rather than a sentence --- the step name \emph{is} the
reason, and the step does not have to remember to report anything.

A \textbf{skip} is about one value. The reading stays in the run with everything else it
carries, and one field is set to \texttt{None}.

Setting the field to \texttt{None} rather than deleting the reading matters, for three reasons.
A \texttt{TankReading} carries conductivity, pH, temperature and weight together, and may carry
the lab sample drawn at that minute --- throwing the row away over one bad pH would throw away
good values and possibly a lab result that cost money to get. \texttt{None} is already a legal
state for every one of those fields, because the loader has to represent a blank cell in the
CSV the same way, so nothing downstream needs new handling for a cleared value. And
\texttt{Check.find\_judgeable\_readings} skips \texttt{None} values by definition, so a cleared
value is invisible to every later check and to every later median without any bookkeeping: it
is simply a reading that was never taken.

\subsection{What it does to the real runs}

Running \texttt{uv run limodel analyze} on the project data gives 32 journal lines before the
results table: 27 skipped readings and 5 dropped runs. All 32 come from this stage --- the
calculating and the fitting share the same journal and wrote nothing to it on this data. The
loader hands over 21 runs carrying 1049 tank readings and 891 power-log rows, and the stage
returns 16 runs --- the line the report prints is \texttt{using 16 of 21 runs}.

Step 1 drops nothing, because all 21 runs are marked for use. All five drops come from step 3:
runs \texttt{AU}, \texttt{AV}, \texttt{AW}, \texttt{AX} and \texttt{AY} do not have both metals
reported at both ends of both tanks. Step 2 drops nothing by construction and accounts for all
27 skips --- 13 tank values and 14 power-log values, out of 1940 readings.

Table~\ref{tab:skips-by-check} is which check found them.

\begin{table}[htbp]
  \centering
  \caption{The 27 skipped values by the check that rejected them, and by the field cleared.}
  \label{tab:skips-by-check}
  \begin{tabular}{@{}lrlr@{}}
    \toprule
    Check & Values & Field cleared & Values \\
    \midrule
    \texttt{PairedValuesTooFarApart}     & 9 & \texttt{meter\_voltage\_v}   & 8 \\
    \texttt{JumpThatComesBack}           & 7 & \texttt{conductivity\_us\_cm} & 7 \\
    \texttt{ValueFarFromTheMedian}       & 6 & \texttt{ph}                  & 3 \\
    \texttt{PairedValuesOutOfProportion} & 2 & \texttt{weight\_g}           & 3 \\
    \texttt{ValueOutsideTheRange}        & 2 & \texttt{supply\_current\_a}  & 3 \\
    \texttt{ValueMovingTooFast}          & 1 & \texttt{meter\_current\_ma}  & 2 \\
                                         &   & \texttt{supply\_voltage\_v}  & 1 \\
    \midrule
    Total                                & 27 &                             & 27 \\
    \bottomrule
  \end{tabular}
\end{table}

The journal lines themselves say what happened in plain words. Three of the real ones:

\begin{quote}\footnotesize
\begin{verbatim}
skipped  reading  AJ:conc:0
    weight 209.7 g is 1802.15 g away from the usual 2011.85 g
skipped  reading  AS:dil:255
    conductivity jumped to 173997 and the next reading
    put 100% of it back
skipped  reading  AG:225
    the two voltage readings differ by 90.06 V, and the
    supply's 100.06 V is unlike its own column
\end{verbatim}
\end{quote}

Read them as three different shapes of mistake. The first is a decimal point in the wrong
place: $209.7\unit{g}$ in a tank whose next recorded weight is $2098.9\unit{g}$ and whose usual
weight is about $2012\unit{g}$. The second is a digit too many, and it is only findable because
the series carries straight on afterwards: $17549.7$, then $173997.1$, then $17307.2$. Nothing
here has to know what the right value was, only that the series left its track for exactly one
reading. The third is a supply voltage of $100.06\unit{V}$ where the stack was running near
$11\unit{V}$, and the multimeter stayed on its own track, so the supply column is the one that
slipped.

The id tells you where to look in the sheet. An id of \texttt{AJ:conc:0} means run
\texttt{AJ}, concentrate tank, minute 0. An id of \texttt{AG:225} has no tank in it because it
is a power-log row: run \texttt{AG}, minute 225. Every line is a cell to go and check, which is
the whole purpose of keeping the journal rather than cleaning silently.
